\section{Cote 44, arithmetic slopes, and integral Lefschetz operators}
\label{sec:arithmetic-sl2}

The conclusions in this section have three distinct statuses.  The literal
height comparison in cote 44, Theorem 5, has a geometric counterexample under
all its explicitly transcribed hypotheses.  Several elementary integral
Lefschetz statements admit complete proofs and exact matrix checks.  A
comparison between an actual theta-lift chain and those Lefschetz lattices
remains an additional mathematical assertion.  The accepted RH/GRH input is
not questioned or audited here.

\subsection{AR-01: a counterexample to the displayed higher-cycle height formula}
\label{ar:counterexample}

\paragraph{Statement.}
Over an algebraically closed field $k$ of characteristic zero there is a
smooth projective fourfold $X$ with a Lefschetz pencil of hyperplane sections,
all of whose fibres are geometrically integral, and an integral class
$Z\in\mathrm{CH}^2(X)$ such that
\[
 HZ=0,\qquad Z|_Y=0\text{ in }\mathrm{CH}^2(Y),\qquad \deg Z^2=3,
\]
where $Y$ is the generic member and $H$ is the hyperplane class.  Consequently
the displayed assertion of cote 44, scan 004, Theorem 5,
\[
 \langle Z|_Y,Z|_Y\rangle=(-1)^m\deg Z^2
\]
is false for $m=2$ under its explicitly written assumptions.  This also
refutes H9(a)'s assertion that integral fibres alone remove the need for
vertical corrections in higher codimension.

\begin{proof}
Let $X=\mathbb P^2_k\times\mathbb P^2_k$, with hyperplane classes $a,b$,
and embed it by the Segre polarization $H=a+b$.  Its Chow ring is
\[
 \mathrm{CH}^*(X)=\mathbb Z[a,b]/(a^3,b^3),\qquad \deg(a^2b^2)=1.
\]
For $Z=a^2-ab+b^2$ one computes
\[
 HZ=a^3+b^3=0,\qquad
 Z^2=a^4-2a^3b+3a^2b^2-2ab^3+b^4=3a^2b^2.
\]
Consider the pencil
\[
 \lambda(x_0y_0+x_1y_1+x_2y_2)
       +\mu(x_1y_1+2x_2y_2)=0.
\]
Its coefficient matrix is
$\operatorname{diag}(\lambda,\lambda+\mu,\lambda+2\mu)$.  It has rank
three except at $[0:1],[-1:1],[-2:1]$, where it has rank two.  A bilinear
polynomial of rank at least two is geometrically irreducible: a nontrivial
factorization of its bidegree $(1,1)$ must have factors of bidegrees $(1,0)$
and $(0,1)$, which would force rank one.  Every member is therefore
geometrically integral.

A rank-three member is smooth.  A rank-two diagonal member has exactly one
singular point, $(e_i,e_i)$, corresponding to the vanishing diagonal entry.
On the chart $x_i=y_i=1$, its equation is the sum of two nonzero multiples
of $x_jy_j$, $j\ne i$.  This is a nondegenerate quadratic form in four
variables, hence an ordinary double point.  None of the three singular
points lies on the base locus: at it $\sum x_jy_j=1$.

The two defining divisors have no common component, so the base is a
codimension-two complete intersection.  If its gradients were dependent
at a base point, that point would be singular on a member of the pencil,
contrary to the preceding description.  The base is smooth.  Its blow-up
therefore gives a Lefschetz pencil.  Each fibre is isomorphic to the
corresponding member, because the base is a Cartier divisor in that member.

Put $K=k(\mathbb P^1)$.  The generic coefficient matrix is invertible, so
a $K$-linear change of $y$ coordinates identifies $Y$ with
$\sum x_i y_i=0$.  Projection to $\mathbb P^2_x$ makes $Y$ the bundle of
lines in
\[
 E=\ker\bigl(\mathcal O^3\longrightarrow\mathcal O(1)\bigr),\qquad
 c(E)=1-a+a^2.
\]
The class $b$ is the first Chern class of the dual tautological line.
The projective bundle relation is
\[
 b^2+c_1(E)b+c_2(E)=b^2-ab+a^2=0.
\]
For clarity, this sign also follows directly from the tautological exact
sequence: its two line factors have first Chern classes $-b$ and $-a+b$,
so $a^2=c_2(E)=(-b)(-a+b)$.  Thus $Z|_Y=0$ in the Chow group.  Every
bilinear height pairing on the restricted classes sends $(0,0)$ to zero,
whereas the displayed formula would give $0=3$.
\end{proof}

The Chow calculation uses the usual projective bundle formula
\cite[Theorem 3.3(b)]{Fulton}; the particular relation is derived above.
This is a counterexample to the raw identity, not to the Hodge sign:
$Z$ is primitive and its square has the correct positive sign.  The generic
fibre is cellular, so its algebraic intermediate Jacobian in this degree
is zero.  In particular, an injective isometric realization of all primitive
middle cycles through that Jacobian is impossible.

The source's margin ends with an ellipsis after the statement that the
fibres are geometrically integral.  No claim is made that an untranscribed
admissibility condition could not have been intended.  In codimension
$m\ge2$, an integral fibre can still support codimension-$(m-1)$ cycles,
which become vertical codimension-$m$ cycles of the total space.  Fibre
integrality removes additional vertical divisors on a surface fibration;
it does not remove these higher-cycle correction groups.

\subsection{AR-02: the sign is absent from slope and norm data}
\label{ar:slope-sign}

\paragraph{Statement.}
A Hermitian lattice, an injective map to its Hermitian dual, all its local
operator norms, and Bost's slope inequality do not determine the sign of
the associated symmetric bilinear form.

\begin{proof}
Take $E=\mathbb Z$ with its standard metric and $E^\vee$ with the dual
metric.  The maps $u_+(1)=1$ and $u_-(1)=-1$ are integral isomorphisms of
norm one at every finite and infinite place.  Their degrees, maximal slopes
and map heights coincide and are zero; both satisfy the slope inequality
with equality.  Their symmetric forms are $xy$ and $-xy$, with opposite
signatures.  Multiplication of a map by $-1$ preserves all these norm data.
\end{proof}

Bost's precise statement for Hermitian bundles $E,F$ over $\mathcal O_K$
and an injective $K$-linear map $\varphi$ is
\[
 \widehat{\deg}_{n}(E)
 \leq\operatorname{rk}(E)
       \bigl(\widehat\mu_{\max}(F)+h(E,F,\varphi)\bigr).
\]
The locator is \cite[\S4.1.4, Proposition 4.5, equation (4.11)]{Bost};
the filtered version is Proposition 4.6, equation (4.14).  Their complete
printed proofs were read: the first uses the saturated image, determinant
norm bounds and maximal slope; the second applies it to filtered quotients
and sums degrees.  Basic arithmetic-degree properties are cited background.

\paragraph{The retrieved step 13 declaration.}
The actual files in \path{Formalism/LatticeSlopes} sharpen this scope.
In \texttt{Statement.lean}, \texttt{Formalism.F.SlopeInequality} concerns
discrete $\mathbb Z$-lattices in finite-dimensional positive Euclidean
spaces.  If $L\ne0$ and $\varphi:L\to M$ is an injective integral map
with $\|\varphi x\|\le C\|x\|$, it asserts
\[
 \mu(L)\le\mu_{\max}(M)+\log C,\qquad
 \deg(L)=-\log\operatorname{covol}(L),\quad
 \mu(L)=\deg(L)/\operatorname{rk}(L).
\]
The hypotheses imply $C>0$.  The complete proof body
\texttt{slopeInequality\_of} in
\path{LatticeSlopesProofs/SlopeInequality.lean} was read.  It proves
\[
 \operatorname{covol}(\varphi L)
 \le C^{\operatorname{rk}(L)}\operatorname{covol}(L)
\]
using a Gram determinant, and applies the maximum-slope bound to the image
lattice.  A positive inner product is an input, not a conclusion about an
algebraic-cycle intersection form.  This was source inspection, not a Lean
kernel run.

The separate file
\path{ReciprocityTheory/ReductiveGroups/ArithmeticGeometry/}
\path{ArakelovGeometry/MoriwakiSt.lean} defines
\texttt{MoriwakiHSStatement} and \texttt{MoriwakiBigStatement}.  These
are arithmetic Hilbert--Samuel asymptotics and a volume bound under
semipositivity and nefness hypotheses; the asymptotic statement also assumes
nonnegative heights.  Neither statement asserts a primitive intersection
signature.  Their full upstream proofs were not inspected.

\subsection{AR-03: metric dependence of an arithmetic self-intersection}
\label{ar:metric}

Let $X=\mathbb P^1_{\mathbb Z}$ and let $\overline L_0$ be
$\mathcal O(1)$ with a smooth conjugation-invariant metric.  Scale the
squared norm by $e^{-t}$ and call the result $\overline L_t$.  With
$x=\widehat c_1(\overline L_0)$ and the convention
$\widehat{\deg}a(\eta)=\frac12\int\eta$, one has
\[
 \widehat{\deg}\,\widehat c_1(\overline L_t)^2
       =\widehat{\deg}(x^2)+t.
\]
\begin{proof}
The representative $-\log\|s\|^2$ for the Green function gains the
constant $t$, so $\widehat c_1(\overline L_t)=x+a(t)$.  The arithmetic
product gives $x\,a(t)=a(t\,c_1(\overline L_0))$ and $a(t)^2=0$,
since $dd^ct=0$.  Expanding the square and using
$\int_{\mathbb P^1(\mathbb C)}c_1(\overline L_0)=1$ proves the formula.
\end{proof}
The product conventions are those of
\cite[\S4.2.11 and \S4.3, especially Theorem 4.3.2]{GilletSoule}.
Thus an arbitrary sign of this arithmetic self-intersection is compatible
with identical finite fibres and Frobenius data.  This does not contradict
a properly primitive arithmetic Hodge-index theorem with its ampleness and
metric hypotheses.

\subsection{AR-04: known cases and the remaining arithmetic comparison}
\label{ar:known}

Cote 44, Theorem 4, concerns the canonical self-duality of $P^m(X)$ for
$\dim X=2m-1$ over an algebraically closed field; the field hypothesis is
explicit in scan 008.  It is not by itself an Arakelov statement over
$\mathcal O_K$.  Theorem 5 concerns function-field heights in a pencil.
The primitive-decomposition margin in scan 008 must remain part of any
reconstruction of the sign convention.

The following are cited background, rather than new consequences of the
accepted RH/GRH theorem.
\begin{itemize}
\item For $m=1$, a curve Jacobian has its classical principal polarization
\cite[Theorem 6.6 and Summary 6.11]{MilneJacobian}.
For a cellular generic fibre in AR-01 the higher algebraic Jacobian is zero,
so its polarization assertion is vacuous.
\item Moriwaki's Theorem A(2) proves, for a regular projective flat
arithmetic variety of relative dimension $d\ge1$ with arithmetically ample
$\overline H$, that nonzero $x\in\widehat{\mathrm{CH}}^1(X)_{\mathbb R}$
satisfying $\widehat c_1(\overline H)^d x=0$ obeys
\[
 \widehat{\deg}\bigl(x^2\widehat c_1(\overline H)^{d-1}\bigr)<0.
\]
Theorem B gives the weak inequality for generic-degree-zero divisors and
characterizes equality by a positive multiple pulled back from
$\operatorname{Spec}\mathcal O_K$ \cite[Theorems A and B]{Moriwaki}.
The exact hypotheses and statements were checked; the complete proof was
not re-proved.
\item K\"unnemann proves the arithmetic standard conjectures for projective
space with the specified canonical Fubini--Study metric
\cite[Theorem 4.1]{Kunnemann}.  His discussion of abelian schemes keeps
the positivity of heights of higher homologically trivial cycles as an
additional conjectural input.  It must not be reported as a theorem for all
higher cycles on abelian varieties.
\item In surface pencils, divisor-height identities include vertical
component corrections.  For the elliptic case, the exact primary locators
are \cite[Lemma 8.1, Theorems 8.4 and 8.6, equations (8.8), (8.11)--(8.14)]{Shioda}.
The correction coefficients come from the negative inverse component
intersection matrix.  These statements were checked at the primary-source
text level, not re-proved in full.
\end{itemize}

No new higher polarization or arithmetic Hodge-index case follows from the
two bare steps 5 and 13.  This is a statement about the missing derivation,
not a claim that every possible stronger comparison is false.  A corrected
higher version needs admissible cycle extensions, vertical correction
groups, and an identity with the particular canonical self-duality and its
primitive signs.  AR-01 proves that the uncorrected universal identity
cannot supply that identity.

\subsection{AR-05: a precise positive consequence of divisibility and slopes}
\label{ar:divisibility-slope}

Let $L,M$ be finite-rank discrete lattices in positive Euclidean spaces,
with $L\ne0$.  If $\varphi:L\to M$ is injective, $\|\varphi\|\le C$
with $C>0$, and
$\varphi L\subset dM$ for a positive integer $d$, then
\[
 \log d\le\mu_{\max}(M)-\mu(L)+\log C.
\]
\begin{proof}
The map $\varphi/d$ is integral and injective, with norm at most $C/d$.
Apply the inspected slope inequality and rearrange.
\end{proof}
If $L$ and $M$ have the same rank $r$, one has the sharper estimate
\[
 \log[M:\varphi L]\le r\log C+\deg(M)-\deg(L).
\]
\begin{proof}
The extension to real spans has absolute determinant equal to the product
of its singular values, hence at most $C^r$.  Therefore
$\operatorname{covol}(\varphi L)\le C^r\operatorname{covol}(L)$.
Use $\operatorname{covol}(\varphi L)=[M:\varphi L]\operatorname{covol}(M)$
and take logarithms.  If $\varphi L\subset dM$, then
$[M:\varphi L]\ge d^r$, giving the corresponding bound with $\mu(M)$.
\end{proof}
This explains how an arithmetic divisibility lower bound can compete with
an archimedean norm upper bound.  It still does not determine a signature:
the maps $\pm d\,\mathrm{id}$ on $\mathbb Z^r$ have the same exact
index $d^r$, norm $d$ and slopes, with equality in these estimates, but
opposite associated symmetric forms.

\subsection{SL-01: integral blocks for divided-power Lefschetz maps}
\label{sl:blocks}

Let $M=\bigoplus_{a=1}^n(\mathbb Ze_a\oplus\mathbb Zf_a)$,
$p_a=e_a\wedge f_a$, $\xi=\sum_a p_a$, and $L(u)=\xi\wedge u$.
The divided power $L^{(m)}=L^m/m!$ is integral because
\[
 \xi^m/m!=\sum_{|R|=m}p_R.
\]
Fix $s$ pairs on which a basis wedge uses exactly one generator, and choose
which generator in each.  Set $v=n-s$.  The degree-$j=s+2t$ block is
indexed by $t$-subsets $T$ of the remaining pairs.  The degree-$j+2m$ block
is indexed by $(t+m)$-subsets $K$, and its matrix is the transpose of the
inclusion matrix $W_{t,t+m}(v)$.  Each such block occurs with multiplicity
$\binom ns2^s$.

\begin{proof}
Expanding $\xi^m/m!$, terms meeting an occupied pair vanish.  Each empty
$m$-set contributes the basis vector indexed by $T\cup R$, with coefficient
one.  A complete pair has even exterior degree, so moving it past any
other complete pair or singleton introduces no sign.  The singleton set,
its choices, and $T$ uniquely describe every exterior basis vector.
Different such singleton data never mix, proving the integral direct sum.
\end{proof}

For $t\le k\le v-t$, Wilson's diagonal entries for $W_{t,k}(v)$ are
\[
 \binom{k-r}{t-r}\quad\hbox{with multiplicity}\quad
 \binom vr-\binom v{r-1},\qquad 0\le r\le t.
\]
The exact verified primary locator is
\cite[Corollary 3, p.~10]{Ghorbani}, which proves the diagonal form of
\cite{Wilson}.  The original Wilson proof was not opened in this run.
The range condition is necessary; H8's unrestricted version omits it.
For hard Lefschetz, $j=i\le n$ and $m=n-i$, every block has $k=v-t$,
so the condition holds.  A diagonal form is not necessarily a Smith form:
entries $2,3$ normalize to invariant factors $1,6$.

\subsection{SL-02: exact integral hard-Lefschetz consequences}
\label{sl:integral}

For $0\le i\le n$, the map
$L^{(n-i)}:\bigwedge^iM\to\bigwedge^{2n-i}M$ is a
$\mathbb Z$-isomorphism exactly for $i=0,1,n$, with repetitions understood.
For $i=2\le n$, its cokernel is $\mathbb Z/(n-1)$.  Over a commutative
ring $A$, all the \emph{ordinary} hard-Lefschetz maps $L^{n-i}$ are
isomorphisms if and only if $n!$ is a unit in $A$.

\begin{proof}
For $i=0,1$, all blocks have $t=0$, and are $1\times1$ unit matrices.
For $i=n$ the operator is the identity.  If $2\le i<n$, choose
$s=i-2,t=1,v=n-i+2=m+2$.  The block $W_{1,v-1}(v)$ has a diagonal entry
$v-1=m+1>1$, so its cokernel is nonzero.  If $i=2$, this is the only
nonunit block, with $v=n$.  Complementing its column indices gives
$J_n-I_n$, whose Smith form is $(1,\ldots,1,n-1)$.

Necessity of the final assertion follows from $L^n$ on degree zero, which
is multiplication by $n!$.  For sufficiency, the corresponding diagonal
entry of an ordinary map is
\[
 m!\binom{t+m-r}{t-r}=\frac{(t+m-r)!}{(t-r)!}.
\]
It is a product of integers between $1$ and $n$.  If $n!$ is a unit,
each factor is a unit, and the integral diagonal equivalences remain
invertible after base change.
\end{proof}

For a principally polarized complex abelian variety of dimension $g$,
the integral exterior description of cohomology therefore gives cokernel
$\mathbb Z/(g-1)$ for cup product by $\theta^{g-2}/(g-2)!$ on degree
two.  This is a cohomological assertion.  It does not assert algebraicity
of arbitrary divided powers in the Chow ring.

\subsection{SL-03: exact matrices for small powers}
\label{sl:computations}

The reproducible script \texttt{wilson\_checks.py} constructs and checks
16 incidence matrices, aggregates all 27 hard-Lefschetz cases $1\le n\le6$,
and independently constructs the full signed exterior matrices in all 14
cases $n\le4$.  All comparisons pass in exact integer arithmetic.  Smith
forms of diagonal matrices are computed by sorting the valuations separately
at each prime; the direct exterior matrices are checked independently.
The complete results are in the three \texttt{wilson\_*\_checks.tsv}
tables.  Here $1^r$ means $r$ unit invariant factors.

\begin{center}
\begin{tabular}{rrrll}
\toprule
$n$&$i$&$m=n-i$&Smith form of $L^m/m!$&Smith form of $L^m$\\
\midrule
2&0&2&$1$&$2$\\
3&2&1&$1^{14},2$&$1^{14},2$\\
4&2&2&$1^{27},3$&$2^{27},6$\\
4&3&1&$1^{48},2^8$&$1^{48},2^8$\\
5&2&3&$1^{44},4$&$6^{44},24$\\
5&3&2&$1^{110},3^{10}$&$2^{110},6^{10}$\\
5&4&1&$1^{166},2^{43},6$&$1^{166},2^{43},6$\\
6&3&3&$1^{208},4^{12}$&$6^{208},24^{12}$\\
6&4&2&$1^{429},3^{65},6$&$2^{429},6^{65},12$\\
6&5&1&$1^{584},2^{196},6^{12}$&$1^{584},2^{196},6^{12}$\\
\bottomrule
\end{tabular}
\end{center}

For example the degree-two block at $n=4,m=2$ is $J_4-I_4$; the other
24 degree-two basis directions give unit blocks.  Its Smith form is
$1^{27},3$.  These are exact computations in the specified exterior model.
The integer counting theta lifts in step 10 has not thereby been identified
with the power of this operator.

\subsection{SL-04: ordinary contraction and normalized lowering differ}
\label{sl:contraction}

Set $D=\sum_a\iota_{f_a^\vee}\iota_{e_a^\vee}$ and let $H$ act by
$n-j$ on degree $j$.  Then
\[
 [D,L]=H,\qquad[H,L]=-2L,\qquad[H,D]=2D.
\]
For a primitive vector $v$ of degree $i\le n$, meaning $Dv=0$, put
$d=n-i$.  For $0\le r\le d$,
\[
 DL^rv=r(d-r+1)L^{r-1}v,\qquad
 D^rL^rv=r!\frac{d!}{(d-r)!}v.
\]
\begin{proof}
The commutator is checked on the basis $1,e_a,f_a,e_a\wedge f_a$ of
each exterior factor and then summed.  The other two relations follow from
degrees.  For $r=1$, $DLv=Hv=dv$.  If the first formula holds for $r$, then
\[
 DL^{r+1}v=(LD+H)L^rv
 =\bigl(r(d-r+1)+d-2r\bigr)L^rv
 =(r+1)(d-r)L^rv.
\]
Iteration gives the product in the second formula.
\end{proof}
The normalized inverse-string operator satisfies
$\Lambda_N L^rv=L^{r-1}v$; it differs from $D$ on this step by the
factor $r(d-r+1)$.  The scalar $(r!)^2$ in H8 is only the endpoint case
$d=r$.

For $n=2$, put $\Omega=e_1\wedge f_1\wedge e_2\wedge f_2$.  Then
\[
 L^2(1)=2\Omega,\qquad D^2\Omega=2,\qquad D^2L^2(1)=4.
\]
Thus the ordinary-contraction definition in cote 16, scan 014, cannot
simultaneously satisfy the inverse identity on scan 016.  This flags a
normalization, coefficient-ring or revision-layer issue in the supplied
text; it is not a claim about unseen handwriting.  A correct normalized
universal construction is given in SL-08 below.

\subsection{SL-05: Arthur blocks do not determine integral indices}
\label{sl:arthur-lattice}

An Arthur block $\mu[t]$ has the form
$\mu\boxtimes\operatorname{Sym}^{t-1}(\mathbb C^2)$.  Its $\mathrm{SL}_2$
weights are $t-1,t-3,\ldots,1-t$, and its associated unramified exponents
are their halves \cite[\S2.3, equation (2.3.1)]{Mok}.  The block is not
itself an integral motive with cup-product correspondences.

\paragraph{Integral counterexample, including $m=1,2,3$.}
In $V_d=\operatorname{Sym}^d(\mathbb Q^2)$, consider the monomial and
binomial lattices
\[
 \Gamma_0=\bigoplus_{r=0}^d\mathbb Zx^{d-r}y^r,\qquad
 \Gamma_1=\bigoplus_{r=0}^d
             \mathbb Z\binom drx^{d-r}y^r.
\]
Both are stable under $\mathrm{SL}_{2,\mathbb Z}$ and give the identical
rational representation.  For $d=m+1$ the divided operator
$e^{(m)}=(x\partial_y)^m/m!$ from the line of $y^d$ to the line of
$x^my$ has Smith invariant $m+1$ on $\Gamma_0$ and $1$ on $\Gamma_1$.

\begin{proof}
The monomial lattice is stable by polynomial substitution.  For the second
lattice, the tensors invariant under permutations in $(\mathbb Z^2)^{\otimes d}$
have the orbit sums with $r$ copies of $y$ as a basis.  Symmetric
multiplication sends these to $\binom drx^{d-r}y^r$ and is equivariant.
Its image is therefore a stable lattice.  Direct differentiation gives
$e^my^{m+1}=(m+1)!x^my$.  The target basis vector in $\Gamma_1$ is
$(m+1)x^my$, so the ordinary coefficients are $(m+1)!$ and $m!$;
division by $m!$ gives $m+1$ and $1$.
\end{proof}

\begin{center}
\begin{tabular}{rcll}
\toprule
$m$&Rational representation&Ordinary indices $(\Gamma_0,\Gamma_1)$
&Divided indices $(\Gamma_0,\Gamma_1)$\\
\midrule
1&$\operatorname{Sym}^2$&$(2,1)$&$(2,1)$\\
2&$\operatorname{Sym}^3$&$(6,2)$&$(3,1)$\\
3&$\operatorname{Sym}^4$&$(24,6)$&$(4,1)$\\
\bottomrule
\end{tabular}
\end{center}
The script also constructs these raising matrices and verifies the table;
the output is \texttt{sl2\_lattice\_checks.tsv}.  This is a fully specified
counterexample to a lattice-free inference.  It is not a numerical claim
about the actual theta lift.

A theta chain changes the automorphic group and adjoins blocks, whereas
cup product moves inside a fixed graded realization.  Identifying the
number of lifts with an Arthur-block length or a Lefschetz exponent requires
additional data.  An index of the associative order generated by operators
also needs a specified ambient integral order; it is not determined by the
Smith form of one operator.

\subsection{SL-06: the actual local divisibility statement and its parameters}
\label{sl:actual-code}

The retrieved \path{Formalism/LocalToGlobal/Divisibility/Inputs.lean}
defines \texttt{ColChain} using a length \texttt{mm}, varying dimensions
$c(j)$, rectangular matrices $K_j$ of size $c(j+1)\times c(j)$, and a
final period functional.  Its field \texttt{fact} identifies the original
period with that functional on the product of the matrices.  It includes
no Lefschetz grading, lowering operator or $\mathfrak{sl}_2$ relation.

Here is the exact scalar content of the inspected local estimate.  Let
$w\mid\ell$ be a place of a number field $K$, with ramification index
$e_w$, normalized by
$|\iota(a)|_\ell=|\ell|_\ell^{v_w(a)/e_w}$.  Suppose $J$ is a finite
nonempty set with a distinguished $j_1$, and a period has a finite expansion
\[
 \iota(P)=\sum_\pi x_\pi y_\pi r_\pi\gamma_{P,\pi}^{-1}
        \prod_{j\in J}
          (\gamma_{\pi,j}^{-1}\varepsilon_{\pi,j}a_{\pi,j}).
\]
Assume that the norms of $x_\pi,y_\pi,r_\pi,\varepsilon_{\pi,j}$ are
at most one; $D$ is a nonzero natural number; all denominators
$\gamma_{\pi,j},\gamma_{P,\pi}$ divide $D$;
$\gamma_{\pi,j_1}=1$; and $|a_{\pi,j}|\le|\ell|^{s_j}$.  Then
\[
 P=0\quad\hbox{or}\quad
 v_w(P)\ge e_w\sum_{j\in J}(s_j-v_\ell(D)).
\]
\begin{proof}
Each inverse denominator costs at most $|\ell|^{-v_\ell(D)}$ in norm.
There are at most $|J|-1$ such costs in the product because the distinguished
denominator is one, and at most one further cost from $\gamma_{P,\pi}$.
Every summand is consequently bounded by
$|\ell|^{\sum_js_j-|J|v_\ell(D)}$.  The ultrametric inequality gives
the same bound for the sum.  If $P\ne0$, substituting the valuation
normalization and using $0<|\ell|<1$ gives the claimed inequality.
\end{proof}

The exact locators within the same directory are
\path{JetBounds/IsExpansion.lean}, declarations
\texttt{PathData}, \texttt{PathData.val}, \texttt{StabDvd},
\texttt{IsExpansion}; \path{JetBounds/Statement.lean}, declarations
\texttt{PathBound}, \texttt{LocalValuationBoundAt},
\texttt{PlaceInputAtEntry}; and
\path{JetBounds/JetBoundsProofs/ValuationOfEntryBounds.lean}, theorem
\texttt{valuation\_of\_entry\_bounds}.  That complete proof and the
complete proofs of \texttt{path\_height\_bound},
\texttt{embedding\_lower\_bound} and
\texttt{norm\_inv\_natCast\_le\_of\_dvd} in
\texttt{NormInvNatCastLeOfDvd.lean} were read.  They implement the
argument above.  The upstream entry estimates and full global application
were not re-proved, and no kernel execution is claimed.

For $|J|=1,2,3$, the bounds are respectively
\[
 e_w(s_1-v_\ell D),\quad
 e_w(s_1+s_2-2v_\ell D),\quad
 e_w(s_1+s_2+s_3-3v_\ell D).
\]
These specialize the actual statement while displaying its parameters.
They are lower bounds, not exact valuations or Smith factors.  Integer-valued
$v_w$ permits replacing the right side by its ceiling.

\paragraph{Actual small-parameter data.}
The arithmetic family in
\path{Formalism/LocalRingCardinality/Params.lean}, structure
\texttt{Ranks}, assumes $m$ even and $m\ge2$, as well as
$m+3\le k$, $k+m+4\le t$, $3m<t$, and
$m\le2\operatorname{ord}\le k$.  Thus the $m=1,3$ examples above are
universal algebraic models, not instances of this parameter family.

The concrete \texttt{numEx} in \texttt{NumbersLemmas.lean} has
$m=2,k=8,t=15,nV=24$, host $q=11,d=7$, and distinguished primes
$3697,7393$.  Its ranks are $N(j)=8+16j$, so the two steps have ranks
$8\to24\to40$.  The exact underlying declarations are
\texttt{R0} in \path{SpanPowerMembership/ExampleRanksRankComplement.lean}
and \texttt{hcV}, \texttt{hcB}, \texttt{hcE} in
\path{SpanPowerMembership/ExampleWeightMiddleEntries.lean}, both under
\path{Formalism/LocalRingCardinality}.  They give
\[
 v=(24,20,16,12,8,4,0),\quad b=(91,109),\quad e=(100,118).
\]
The host index set has four elements, so the real-place index set has
eight.  At each index, the rows in \texttt{Numbers.lean} are
\[
\begin{array}{c|l|l}
j&a^+_j&a^-_j\\ \hline
1&(109,106,103,100,97,94,91,1)&(1)\\
2&(126,123,120,117,114,111,108,18,1)&()
\end{array}
\]
and therefore
\[
 (D^+_1,D^+_2)=(5608,6704),\quad
 (D^-_1,D^-_2)=(8,0),
\]
\[
 (F_{G,1},F_{G,2})=(-2808,-3352),\qquad
 (F_{\mathrm{ram},1},F_{\mathrm{ram},2})=(-1400,-1652).
\]
These follow by exact rational substitution, with
$F_{G,j}=-(D^+_j+D^-_j)/2$ and
$F_{\mathrm{ram},j}=-(7/2)\cdot4e_j$.
The first ramified coefficient and the first zero coefficients at $2,11$
also agree with the explicit source declarations
\texttt{ex\_FRamD}, \texttt{ex\_FjP\_d},
\texttt{ex\_FjP\_2}, \texttt{ex\_FjP\_q}.

\paragraph{Exact evaluation of the orbit optimization.}
For constant inputs $D_s=A\ge0$, $p\ge2$ and an orbit length
$d_0\ge2$, the linear program \texttt{lpFeasible} in
\texttt{DataParams\textsubscript{2}.lean} has optimum
$d_0A/(p+1)$, and
\[
 \operatorname{lpF}(p,d_0,A)=\frac{d_0A}{p^2-1}.
\]
\begin{proof}
The variables $x_{s,i}\ge0$, indexed by $s\in\mathbb Z/d_0\mathbb Z$,
vanish unless $i<d_0$ is odd.  The constraints are
\[
 \sum_{i<d_0}p^i x_{s,i}+\sum_{i<d_0}x_{s-i,i}\le A.
\]
Summing over $s$ and translating the second sum gives
\[
 (p+1)\sum_{s,i}x_{s,i}
 \le\sum_{s,i}(p^i+1)x_{s,i}\le d_0A,
\]
since every supported $i$ is at least one.  The point
$x_{s,1}=A/(p+1)$, all other variables zero, satisfies every constraint
with equality and attains this bound.  Substituting the optimum in
$\operatorname{lpF}=p d_0A/(p^2-1)-\sup\operatorname{lpFeasible}$
gives the formula.
\end{proof}
Here $A=7e_j$.  If the inert class has zero Frobenius element, the
exceptional branch of \texttt{FInert} gives the same result directly.
In the remaining inert or signal cases the orbit length is at least two,
and the displayed lemma applies.  Consequently, at any such prime,
\[
 F_{\mathrm{inert},j}(p)=\frac{56e_j}{p^2-1}.
\]

For \texttt{numEx}, the cutoff is
$X''=8\exp(2^{1/300})$, with $21<X''<22$.
For a rational certificate, $(101/100)^{300}>2$ and the Taylor sum
through degree $12$ for $\exp(101/100)$, with the geometric bound on its
remaining tail, is less than $11/4$; also $e>8/3$.
The prime classification is the one in \texttt{HostCls.lean}:
outside $2,7,11$, split primes have $(-7/p)=1$; the others are signal
or inert according as $(-11/p)=1$ or $-1$.
This gives the entire nonzero contribution table for this numerical datum:
\begin{center}
\begin{tabular}{rcrrr}
\toprule
$p$&Class&$F_{1,p}$&$F_{2,p}$&$\sum_jF_{j,p}$\\
\midrule
2&ramified at $2$&0&0&0\\
3&signal&700&826&1526\\
5&signal&$700/3$&$826/3$&$1526/3$\\
7&ramified at $d$&$-1400$&$-1652$&$-3052$\\
11&ramified at $q$&0&0&0\\
13&inert&$100/3$&$118/3$&$218/3$\\
17&inert&$175/9$&$413/18$&$763/18$\\
19&inert&$140/9$&$826/45$&$1526/45$\\
3697,7393&distinguished&$-2808$&$-3352$&$-6160$\\
\bottomrule
\end{tabular}
\end{center}
For each of the last two primes the row applies separately.  In particular,
the actual definition of \texttt{Fo} evaluates to
\[
 12208\sum_{p\in\{3,5,13,17,19\}}\frac{\log p}{p^2-1}
 -3052\log7-6160(\log3697+\log7393).
\]
Its decimal approximation is $-108518.601323$; this is the $F$
contribution alone.  The exact evaluation, the cutoff certificate and
twenty primal/upper-bound certificates for the relevant orbit programs
are reproduced by \texttt{step10\_numeric\_checks.py}.  The parameter
and prime tables are saved separately with hashes of the inspected source
files.  No Lean kernel execution is claimed.

\paragraph{The exact remaining hypotheses.}
The small tuple \texttt{numEx} is a test datum, and cannot satisfy
\texttt{NumFacts} in \texttt{NumbersLink.lean}: its field \texttt{n2}
would require $N(m)^2\le b$, whereas
\[
 N(m)^2=40^2=1600,\qquad b=\exp(2^{1/100})<e^2<9.
\]
This proves incompatibility of that particular tuple with these hypotheses;
it does not refute the conditional theorem.  The same file explicitly
provides the large-scale repair \texttt{NumInst.numBig}, with
$M=200+|\texttt{numEx.logOm}|$ and $x=\exp(100M)$, and constructs
\texttt{numFactsBig : numBig.NumFacts wd0 1 1 5 5 1 53 4}.
Its $v,b,e$ are unchanged; only one $u$ entry and the scale change.
Thus the distinguished and ramified coefficients above and the formula
$56e_j/(p^2-1)$ also hold for this larger example, with further primes
admitted by its larger cutoff.  This is a check of the exact declaration
and its parameters, not a full check of all supporting \texttt{numFactsBig}
proofs.

Neither a numerical datum nor \texttt{NumFacts} selects a nonzero period
and its two integral lattices.  The exact record-level definition in
\texttt{DataParams\textsubscript{2}.lean} is
\[
 \operatorname{sg}(u,w)=\sum_{j=1}^m
   \bigl(\operatorname{gainD}(j,w)+F(j,w)-E(u,j,w)\bigr),
\]
where \texttt{gainD} suppresses the gain at primes in \texttt{drop}.
The global gain term also subtracts \texttt{tau}; it is not part of this
local definition.  One still needs the record linking its $F$ to the
numeric $F_{j,p}$, the gain and error data, the period expansion, entry
estimates, integral row and column lattices, stabilizer denominator, and
the inequality relating those inputs to the requested local threshold.
For comparison, the companion \texttt{ExampleWeightMiddleEntries.lean}
also constructs an \texttt{emptyColumnChain} and
\texttt{trivialPeriodMap} whose period is identically zero.  These examples
are not a supplied nonzero theta-lift matrix with a specified Smith form.

At the level of matrix-chain data and factorization, there is no invariant
depending only on length.  Fix $m\ge1$, set every dimension equal to one,
take all frames and all matrices equal to one except $K_0=[\ell^s]$, and
take the final functional to be multiplication.  The period is
$\ell^sxy$, satisfying the factorization and taking valuation $s$ at
$x=y=1$.  Every $s\ge0$ is possible at fixed length.  This refutes a
length-only conclusion from that bare matrix structure, not from the
additional automorphic hypotheses of a particular lift.

\paragraph{The actual divided-power action.}
In \path{Formalism/GradedWeylDivisibility/StatementA.lean},
\texttt{rowDiv R a b k f} is the coefficient of $T^k$ after the row
transvection $X_{b,s}\mapsto X_{b,s}+TX_{a,s}$.  The Weyl lattice is
generated by words in these operators on a highest-weight polynomial;
the dual Weyl lattice is defined by column homogeneity and unipotent
invariance.  \texttt{StatementB.lean} uses a factorial-weighted Fock
pairing, and \texttt{StatementD.lean}, \texttt{StatementE.lean}
treat transformed integral lattices and duality.  These are concrete
integral general-linear-group modules.  No identification with the
symplectic exterior operator $L$ is present in the inspected definitions.

An equality with Wilson's factors therefore needs a map to a fixed graded
geometric realization, equality of the actual operators with specified
divided powers, and preservation of the specified integral lattices and
denominators.  The numerical input not yet supplied is a comparison gate,
not itself a mathematical counterexample.  A scoped search of the three
divisibility folders, \texttt{PositivityTheory} and
\path{Formalism/HermitianCurvature} found no named arithmetic Hodge-index
comparison; this is not a nonexistence assertion about the entire archive.

\subsection{SL-07: a Frobenius normalization obstruction}
\label{sl:frobenius}

Let $X/\mathbb F_q$ be smooth projective of positive dimension, and let
$\eta$ be an ample hyperplane class defined over $\mathbb F_q$.  Use
cohomological geometric Frobenius $F$ on untwisted $\ell$-adic cohomology,
so $F\eta=q\eta$.  Then for $Lx=\eta\cup x$,
\[
 FL=qLF.
\]
\begin{proof}
Pullback by the $q$-power Frobenius sends $\mathcal O(1)$ to
$\mathcal O(q)$.  Chern-class functoriality gives $F\eta=q\eta$,
and cup-product functoriality proves the identity.  Equivalently, the class
is Galois-invariant in $H^2(\overline X,\mathbb Q_\ell(1))$.
\end{proof}
The conventions and Chern-class facts were checked in
\cite[Theorem 23.3 and \S27, especially Proposition 27.3]{MilneEtale}.

In contrast, the two factors of a local parameter
$\psi:L_F\times\mathrm{SL}_2(\mathbb C)\to{}^LG$ commute.  Thus the
raising element commutes with $\psi(w,1)$.  A linear identification of
that pair with the untwisted pair $(L,F)$ would imply
$(q-1)LF=0$.  Since $F$ is invertible and coefficients have characteristic
zero, it would force $L=0$.  Already $\mathbb P^1$ gives
$F=\operatorname{diag}(1,q)$ and a nonzero map $1\mapsto\eta$.

This only obstructs the unnormalized identification.  On
$\bigoplus_rH^{2r}(\overline X,\mathbb Q_\ell(r))$, and separately on
$\bigoplus_rH^{2r+1}(\overline X,\mathbb Q_\ell(r))$, hyperplane
multiplication commutes with Frobenius.  These integer Tate twists are
also integral over $\mathbb Z_\ell$ when $\ell\ne\operatorname{char}\mathbb F_q$.
One must not call every such normalization nonintegral.  A uniform
weight-zero normalization can additionally require half twists.
Furthermore the associated Langlands parameter is
\[
 \varphi_\psi(w)=\psi\left(w,
 \begin{pmatrix}|w|^{1/2}&0\\0&|w|^{-1/2}\end{pmatrix}\right),
\]
not $\psi(w,1)$, and need not commute with the raising element
\cite[\S2.2, equations (2.2.9)--(2.2.10)]{Mok}.  Proper normalization
removes this particular obstruction, but does not provide the missing
integral geometric comparison.

\subsection{SL-08: the normalized universal ring, with complete proof}
\label{sl:universal}

Define
\[
 \Phi_n^{\mathrm{str}}=\prod_{i=0}^n M_{n-i+1}(\mathbb Z).
\]
Its $i$th factor acts on a string $v_{i,0},\ldots,v_{i,n-i}$ in degrees
$i,i+2,\ldots,2n-i$.  Let $L$ move one step up and $\Lambda_N$ one
step down, both zero outside the string.  Write $E_i^{\beta,\alpha}$
for its matrix units, of degree $2(\beta-\alpha)$, and let $\pi_j$
be the sum of diagonal units in degree $j$.  Then
\[
 E_{i'}^{\beta',\alpha'}E_i^{\beta,\alpha}
   =\delta_{i',i}\delta_{\alpha',\beta}E_i^{\beta',\alpha}.
\]
The ring is generated over $\mathbb Z$ by $L,\Lambda_N$ and is free of
rank
\[
 \sum_{i=0}^n(n-i+1)^2=\frac{(n+1)(n+2)(2n+3)}6.
\]

\begin{proof}
Set $B=1-L\Lambda_N$, the projection to all string bottoms.  For
$0\le d\le n$,
\[
 Q_d=B(\Lambda_N^dL^d-\Lambda_N^{d+1}L^{d+1})B
\]
projects to the bottom of the string of length $d+1$.  Indeed the two
terms on a bottom vector are the indicators for lengths at least $d+1$
and at least $d+2$.  Every matrix unit on this string is
$L^\beta Q_d\Lambda_N^\alpha$.  This proves generation, while the
product of matrix rings gives the asserted free rank.
\end{proof}

The correct string-index range is $0\le\alpha,\beta\le n-i$.
The range $2n-i$ in scan 015 is incompatible with its degrees and its
claimed basis.  The corrected rank is five for $n=1$, agreeing with
scan 017.

\paragraph{Universal property.}
Let $A$ be a commutative ring and $M=\bigoplus_{j=0}^{2n}M^j$ a graded
$A$-module, with a degree-two $A$-linear operator $L$.  Suppose the ordinary
maps
\[
 L^{n-i}:M^i\longrightarrow M^{2n-i}
\]
are isomorphisms for $0\le i\le n$.  Then there is a unique graded
$A$-linear representation of $A\otimes\Phi_n^{\mathrm{str}}$ carrying
$L$ and every degree projector to the prescribed operators.  Its lowering
operator is the normalized inverse-string operator.  Conversely, such a
representation implies these hard-Lefschetz isomorphisms.  For a
unital graded $A$-subalgebra $\mathcal E\subset\operatorname{End}_A(M)$
containing $L$, existence of the
representation into $\mathcal E$ is equivalent to the normalized lowering
operator and all degree projectors belonging to $\mathcal E$.

\begin{proof}
For $i\le n$, put
\[
 P^i=\ker(L^{n-i+1}:M^i\to M^{2n-i+2}),
\]
with out-of-range groups zero.  If $i\ge2$, the map
\[
 x\longmapsto
 (L^{n-i+2}:M^{i-2}\to M^{2n-i+2})^{-1}L^{n-i+1}x
\]
is a left inverse of $L:M^{i-2}\to M^i$, with kernel $P^i$.
Therefore $M^i=P^i\oplus LM^{i-2}$; for $i=0,1$ one has $P^i=M^i$.
Induction proves the decomposition in degrees at most $n$.  Transporting
the decomposition of degree $2n-j$ through $L^{j-n}$ proves it in degree
$j>n$.  In total,
\[
 M=\bigoplus_{i=0}^n\bigoplus_{r=0}^{n-i}L^rP^i.
\]
Each $L^r$ is injective on $P^i$ because $L^{n-i}$ is injective there.
Thus matrix units act unambiguously by $L^\alpha v\mapsto L^\beta v$
within each primitive string.  This constructs the representation and
the lowering operator, and immediately gives the converse.

For uniqueness, a representation with the prescribed degree projectors
decomposes $M$ into its matrix-unit summands.  Its bottom summand in the
$i$th factor is exactly $P^i$: $L^{n-i+1}$ kills that summand and is
injective on every other summand in degree $i$.  The remaining summands
are consequently $L^rP^i$, and $L$ forces every matrix unit between them.
The representation is unique.  Since the matrix units are generated by
$L,\Lambda_N$, the subalgebra assertion follows.
\end{proof}

The degree condition is essential.  For $n=1$, conjugation by $1+tL$ in
the $M_2(\mathbb Z)$ factor fixes $L$ but sends
\[
 \Lambda_N\longmapsto
 \Lambda_N+t(E_{11}-E_{00})-t^2L.
\]
Different integers $t$ give different ungraded maps with the same $L$;
they do not preserve the degree projectors.

The symplectic exterior lattice with ordinary $L$ does not furnish this
unit-step model over $\mathbb Z$ in general.  By SL-02, all its ordinary
hard-Lefschetz maps become isomorphisms after inverting $n!$.  Already for
$n=2$, $L^2(1)=2\Omega$ forces $\Lambda_N\Omega=\xi/2$, which is not
integral in the exterior lattice.  After base change to
$\mathbb Z[1/n!]$, the universal property applies.  This repairs the
normalized calculation; it does not identify this ring with the distinct
integral contraction order $\mathbb Z[L,D]$ in an unspecified ambient
matrix order.

\paragraph{Verification boundary.}
The new elementary arguments above were read in full by the authoring
agent.  AR-01 was independently reconstructed and its entire proof checked
by the openness agent; SL-08 was independently re-proved by the sign-geometry
agent.  Literature statement checks, complete proof checks, exact integer
computations and source-level Lean inspection are recorded separately in
the accompanying event ledger.  No full Lean compilation or proof of the
entire imported dependency graph is claimed.
